% !TeX root = ../exemple-B-tp.tex
% !TeX encoding = UTF-8
%% =====================================================================
%%  Exemple B — Compte rendu de travaux pratiques
%%  Sujet : mesure de l'accélération de la pesanteur avec un pendule simple.
%%  Montre : blocs de TP (objectif, principe, protocole, résultats,
%%  analyse, conclusion), tableaux de mesures (siunitx S, booktabs,
%%  tabularx), figure à sous-figures (schéma TikZ + courbe pgfplots),
%%  exemple de calcul avec unités, propagation d'incertitudes.
%%
%%  ATTENTION : toutes les valeurs numériques ci-dessous sont FICTIVES.
%%  Elles servent uniquement à illustrer la mise en forme.
%% =====================================================================

\begin{note}[Données d'illustration]
  Les mesures de ce compte rendu sont fictives : elles montrent la mise en
  page d'un TP, pas le résultat d'une expérience réelle.
\end{note}

%% ---------------------------------------------------------------------
\begin{objectif}
%% ---------------------------------------------------------------------
Déterminer l'accélération de la pesanteur $g$ à partir de la période d'un
pendule simple, en faisant varier sa longueur $L$, et estimer l'incertitude
associée au résultat. On compare ensuite la valeur obtenue à la valeur
conventionnelle $g_{n}=\qty{9.80665}{\metre\per\second\squared}$.
\end{objectif}

%% ---------------------------------------------------------------------
\begin{principe}
%% ---------------------------------------------------------------------
Un pendule simple de longueur $L$, écarté d'un angle initial $\theta_{0}$,
oscille avec la période
\begin{equation}
  T = 2\pi\sqrt{\frac{L}{g}}\,\Bigl(1+\frac{\theta_{0}^{2}}{16}
      +\bigO{\theta_{0}^{4}}\Bigr),
  \label{eq:periode}
\end{equation}
où la correction entre parenthèses, due à l'amplitude finie, est négligeable
pour $\theta_{0}\lesssim\ang{10}$ (voir \cref{sec:discussion}). Au
premier ordre, \cref{eq:periode} se linéarise en
\begin{equation}
  T^{2} = a\,L,
  \qquad a=\frac{4\pi^{2}}{g},
  \qquad\text{d'où}\qquad
  g=\frac{4\pi^{2}}{a}.
  \label{eq:lineaire}
\end{equation}
Tracer $T^{2}$ en fonction de $L$ et ajuster une droite permet donc
d'obtenir $g$ à partir de la pente $a$, sans connaître exactement
l'origine des longueurs.
\end{principe}

%% ---------------------------------------------------------------------
\begin{protocole}
%% ---------------------------------------------------------------------
\subsection{Matériel}

Le matériel utilisé et la résolution de chaque instrument sont rassemblés
dans le \cref{tab:instruments}.

\begin{table}[htbp]
  \caption{Instruments de mesure et résolutions.}
  \label{tab:instruments}
  \begin{tabularx}{\linewidth}{@{}l Y l@{}}
    \toprule
    Instrument & Grandeur mesurée & Résolution \\
    \midrule
    Règle graduée      & Longueur $L$ (point d'attache $\to$ centre de la masse) & \qty{1}{\milli\metre} \\
    Chronomètre manuel & Durée $t_{10}$ de dix oscillations                      & \qty{0.01}{\second} \\
    Rapporteur         & Angle initial $\theta_{0}$                              & \ang{1} \\
    \bottomrule
  \end{tabularx}
\end{table}

\subsection{Mode opératoire}

\begin{enumerate}
  \item Fixer le fil au support et ajuster la longueur $L$ (de \qty{30}{\centi\metre} à \qty{100}{\centi\metre}, par pas de \qty{10}{\centi\metre}).
  \item Écarter la masse d'un angle $\theta_{0}\approx\ang{10}$ et la lâcher sans vitesse initiale.
  \item Mesurer la durée $t_{10}$ de dix oscillations complètes ; en déduire $T=t_{10}/10$.
  \item Reporter les mesures dans le \cref{tab:mesures} et recommencer pour chaque longueur.
\end{enumerate}
\end{protocole}

%% ---------------------------------------------------------------------
\begin{resultats}
%% ---------------------------------------------------------------------
\subsection{Mesures}

Le \cref{tab:mesures} donne les durées mesurées, avec $T=t_{10}/10$. Les
incertitudes-types ont été estimées comme suit : $u(L)=\qty{0.2}{\centi\metre}$
(lecture de la règle et repérage du centre de la masse\footnote{Le centre de la masse est
repéré au demi-millimètre près ; l'incertitude retenue est volontairement plus large.}) et
$u(t_{10})=\qty{0.15}{\second}$ (temps de réaction de l'opérateur), d'où
$u(T)=u(t_{10})/10=\qty{0.015}{\second}$.

\begin{table}[htbp]
  \caption{Durée $t_{10}$ de dix oscillations en fonction de la longueur $L$
    ($\theta_{0}\approx\ang{10}$).}
  \label{tab:mesures}
  \begin{threeparttable}
    \begin{tabular}{S[table-format=3.1] S[table-format=2.2] S[table-format=1.3]
                    S[table-format=1.3] S[table-format=1.3]}
      \toprule
      {$L$ (\unit{\centi\metre})} & {$t_{10}$ (\unit{\second})} & {$T$ (\unit{\second})\tnote{a}}
        & {$T^{2}$ (\unit{\second\squared})} & {$u(T^{2})$ (\unit{\second\squared})\tnote{b}} \\
      \midrule
       30.0 & 10.98 & 1.098 & 1.206 & 0.033 \\
       40.0 & 12.71 & 1.271 & 1.615 & 0.038 \\
       50.0 & 14.17 & 1.417 & 2.008 & 0.043 \\
       60.0 & 15.52 & 1.552 & 2.409 & 0.047 \\
       70.0 & 16.72 & 1.672 & 2.796 & 0.050 \\
       80.0 & 17.93 & 1.793 & 3.215 & 0.054 \\
       90.0 & 19.12 & 1.912 & 3.656 & 0.057 \\
      100.0 & 20.09 & 2.009 & 4.036 & 0.060 \\
      \bottomrule
    \end{tabular}
    \begin{tablenotes}[flushleft]\footnotesize
      \item[a] $T=t_{10}/10$.
      \item[b] $u(T^{2})=2\,T\,u(T)$, avec $u(T)=\qty{0.015}{\second}$.
    \end{tablenotes}
  \end{threeparttable}
\end{table}

\subsection{Exemple de calcul}

Pour la longueur $L=\qty{100.0}{\centi\metre}$, on mesure
$t_{10}=\qty{20.09}{\second}$, donc
\begin{align}
  T   &= \frac{t_{10}}{10} = \frac{\qty{20.09}{\second}}{10} = \qty{2.009}{\second},
  \label{eq:calcT}\\
  g   &= \frac{4\pi^{2}L}{T^{2}}
       = \frac{4\pi^{2}\times\qty{1.000}{\metre}}{(\qty{2.009}{\second})^{2}}
       = \qty{9.78}{\metre\per\second\squared}.
  \label{eq:calcg}
\end{align}
L'incertitude relative composée vaut
\begin{equation}
  \incertrel{g}=\sqrt{\Bigl(\frac{u(L)}{L}\Bigr)^{2}+\Bigl(\frac{2\,u(T)}{T}\Bigr)^{2}}
  =\sqrt{(\num{0.0020})^{2}+(\num{0.0149})^{2}}=\num{0.015},
  \label{eq:uponctuel}
\end{equation}
soit $g=\qtyerr{9.78}{0.15}{\metre\per\second\squared}$ pour cette seule
mesure.

\subsection{Ajustement linéaire}

Les points $(L,T^{2})$ sont représentés sur la \cref{fig:regression}, avec la
droite $T^{2}=aL+b$ obtenue par moindres carrés ordinaires sur les huit
points. Les résultats de l'ajustement sont réunis dans le \cref{tab:ajustement}.

\begin{figure}[htbp]
  \begin{subfigure}[b]{0.34\textwidth}
    \centering
    \begin{tikzpicture}[>=Stealth, line cap=round, scale=0.78]
      \def\Lr{3.4}
      \coordinate (O) at (0,0);
      \coordinate (V) at (0,-\Lr);
      \coordinate (M) at ({\Lr*sin(24)},{-\Lr*cos(24)});
      \draw[line width=0.8pt, TextColor] (-1.1,0)--(1.1,0);
      \foreach \x in {-1.0,-0.7,...,1.0}
        \draw[line width=0.4pt, MutedColor] (\x,0)--(\x+0.18,0.2);
      \draw[dashed, MutedColor, line width=0.4pt] (O)--(V);
      \draw[line width=0.6pt, TextColor] (O)--(M);
      \fill[DataBlue] (M) circle (0.17);
      \pic[draw=DataBlue, line width=0.5pt, angle radius=1.3cm,
           "$\theta_{0}$", angle eccentricity=1.35] {angle=V--O--M};
      \draw[<->, line width=0.4pt, MutedColor]
        ($(O)+({0.5*cos(24)},{0.5*sin(24)})$) --
        node[midway, fill=white, inner sep=1pt] {\footnotesize$L$}
        ($(M)+({0.5*cos(24)},{0.5*sin(24)})$);
      \draw[->, line width=0.7pt, MutedColor] (-1.7,-1.2)--(-1.7,-2.6)
        node[midway, left] {\footnotesize$\vect{g}$};
    \end{tikzpicture}
    \caption{Schéma du dispositif.}
    \label{fig:schema}
  \end{subfigure}\hfill
  \begin{subfigure}[b]{0.62\textwidth}
    \centering
    \begin{tikzpicture}
      \begin{axis}[
          width=\linewidth, height=6.1cm,
          xlabel={$L$ (\unit{\metre})}, ylabel={$T^{2}$ (\unit{\second\squared})},
          xmin=0.2, xmax=1.1, ymin=0.6, ymax=4.6,
          grid=major, legend pos=north west,
          scaled ticks=false,
          /pgf/number format/use comma,
        ]
        \addplot[only marks, mark=*, mark size=1.6pt, color=DataBlue,
                 error bars/.cd, y dir=both, y explicit,
                 error bar style={line width=0.4pt, color=DataBlue},
                 error mark options={rotate=90, mark size=2pt}]
          table[x=L, y=T2, y error=u, col sep=space] {
            L    T2     u
            0.30 1.206  0.033
            0.40 1.615  0.038
            0.50 2.008  0.043
            0.60 2.409  0.047
            0.70 2.796  0.050
            0.80 3.215  0.054
            0.90 3.656  0.057
            1.00 4.036  0.060
          };
        \addplot[domain=0.22:1.08, samples=2, thick, color=DataRed]
          {4.0503*x - 0.0152};
        \legend{Mesures, Ajustement}
      \end{axis}
    \end{tikzpicture}
    \caption{$T^{2}$ en fonction de $L$ et ajustement linéaire.}
    \label{fig:regression}
  \end{subfigure}
  \caption{Pendule simple : (a) schéma du montage, avec l'angle initial
    $\theta_{0}$ et le champ de pesanteur $\vect{g}$ ; (b) mesures de
    $T^{2}$ avec barres d'incertitude $\pm u(T^{2})$ et droite de régression.}
  \label{fig:pendule}
\end{figure}

\begin{table}[htbp]
  \caption{Résultats de l'ajustement $T^{2}=aL+b$ (huit points, six degrés de liberté).}
  \label{tab:ajustement}
  \begin{tabular}{@{}l l l@{}}
    \toprule
    Grandeur & Expression & Valeur \\
    \midrule
    Pente $a$              & $\dfrac{n\sum L_i T_i^{2}-\sum L_i\sum T_i^{2}}{\Delta}$
                           & \qtyerr{4.050}{0.025}{\second\squared\per\metre} \\[2.4ex]
    Ordonnée $b$           & $\dfrac{\sum T_i^{2}-a\sum L_i}{n}$
                           & \qtyerr{-0.015}{0.017}{\second\squared} \\[2.4ex]
    Écart-type résiduel $s$ & $\sqrt{\sum r_i^{2}/(n-2)}$
                           & \qty{0.016}{\second\squared} \\[1.6ex]
    Accélération $g$       & $4\pi^{2}/a$
                           & \qtyerr{9.75}{0.06}{\metre\per\second\squared} \\
    \bottomrule
  \end{tabular}
\end{table}
\end{resultats}

%% ---------------------------------------------------------------------
\begin{analyse}
%% ---------------------------------------------------------------------
\subsection{Analyse des incertitudes}

Avec $\Delta=n\sum L_i^{2}-\bigl(\sum L_i\bigr)^{2}$ et $r_i=T_i^{2}-aL_i-b$
les résidus, les incertitudes-types de la droite de régression sont
\begin{equation}
  s^{2}=\frac{1}{n-2}\sum_{i=1}^{n}r_i^{2},
  \qquad
  u(a)=s\sqrt{\frac{n}{\Delta}},
  \qquad
  u(b)=s\sqrt{\frac{\sum L_i^{2}}{\Delta}} .
  \label{eq:regression}
\end{equation}
Comme $g=4\pi^{2}/a$, la propagation à une seule variable donne
\begin{equation}
  \incertrel{g}=\frac{u(a)}{a}
  \quad\Longrightarrow\quad
  u(g)=g\,\frac{u(a)}{a}
  =\qty{9.75}{\metre\per\second\squared}\times\frac{\qty{0.025}{\second\squared\per\metre}}{\qty{4.050}{\second\squared\per\metre}}
  =\qty{0.06}{\metre\per\second\squared}.
  \label{eq:ug}
\end{equation}
Le résultat final est donc
\begin{equation}
  g=\qtyerr{9.75}{0.06}{\metre\per\second\squared}.
  \label{eq:gfinal}
\end{equation}

\subsection{Discussion}\label{sec:discussion}

\paragraph{Comparaison avec la valeur de référence.}
L'écart normalisé entre \cref{eq:gfinal} et la valeur conventionnelle $g_{n}$ est
\begin{equation}
  E_{n}=\ecartnorm{g}{g_{n}}\approx
  \frac{\abs{\qty{9.75}{\metre\per\second\squared}-\qty{9.80665}{\metre\per\second\squared}}}
       {\qty{0.06}{\metre\per\second\squared}}\approx\num{1.0}.
  \label{eq:En}
\end{equation}
En prenant $u(g_{n})=0$, ce résultat est compatible avec la valeur de
référence au niveau de confiance usuel ($E_{n}<2$). Le protocole
d'évaluation suit le \emph{Guide pour l'expression de
l'incertitude de mesure}~\cite{jcgm2008} ; voir aussi \cite{taylor1997}.

\paragraph{Cohérence des incertitudes.}
L'écart-type résiduel ($s=\qty{0.016}{\second\squared}$) est environ trois fois
inférieur aux incertitudes propagées $u(T^{2})\in[\num{0.033},\num{0.060}]\,\unit{\second\squared}$
du \cref{tab:mesures} : l'estimation $u(t_{10})=\qty{0.15}{\second}$ du temps
de réaction est donc probablement trop prudente. L'incertitude finale, calculée
à partir de la dispersion des points (\cref{eq:regression}), est plus
réaliste que celle qu'on obtiendrait en propageant directement les
incertitudes de type~B.

\paragraph{Erreurs systématiques.}
La correction d'amplitude de \cref{eq:periode} allonge la période mesurée de
$\theta_{0}^{2}/16\approx\num{0.2}\,\%$ pour $\theta_{0}=\ang{10}$, ce qui
sous-estime $g$ d'environ $\theta_{0}^{2}/8\approx\num{0.4}\,\%$, soit
\qty{0.04}{\metre\per\second\squared}. Cette correction reste inférieure à
$u(g)$ mais deviendrait significative pour une mesure plus précise. Les
autres sources possibles sont la masse non nulle du fil, les dimensions finies
de la masse suspendue et les frottements de l'air.
\end{analyse}

%% ---------------------------------------------------------------------
\begin{conclusion}
%% ---------------------------------------------------------------------
La méthode de linéarisation $T^{2}=(4\pi^{2}/g)\,L$ a fourni
$g=\qtyerr{9.75}{0.06}{\metre\per\second\squared}$, compatible avec la valeur
de référence $g_{n}$ (écart normalisé $\approx1$). L'ajustement sur huit longueurs
réduit l'incertitude par rapport à une mesure ponctuelle
(\cref{eq:uponctuel} : \qty{0.15}{\metre\per\second\squared}). Pour améliorer
la précision, on pourrait chronométrer vingt oscillations, réduire l'amplitude
initiale et appliquer la correction d'amplitude.
\end{conclusion}
